\tzpanchor{1071}
\tzpentryheader{1071}{MAGNETIC FIELD RAMP-UP}

\begin{tzpsnapshot}
\tzpsnaprow{TZPID}{\tzpcode{ID1071}}
\tzpsnaprow{UUID}{\tzpcode{1cf54662-1a55-538c-b31c-15ba9d038b6f}}
\tzpsnaprow{Type}{Field / Transport Statement}
\tzpsnaprow{Scale}{magnetofluid}
\tzpsnaprow{LOC Classification}{Working routing: QC809 magnetohydrodynamics; QC6 fluid and plasma dynamics}
\tzpsnaprow{arXiv Classification}{Primary: physics.plasm-ph; Cross-list: astro-ph.SR, gr-qc}
\end{tzpsnapshot}

\tzpsection{Canonical Statement}
B(t) = int\_0\textasciicircum{}t dot\{B\}\_\{ramp\} dt implies B\_\{target\} approximately 0.1 text\{ Tesla\}

\[
B(t) = \in t_0^t \mathrm{dot}{B} {ramp} dt \implies B {target} \approx 0.1 \text{ Tesla}
\]
\[
- Magnetic field B(r,\theta,\phi) =  abla \times A
\]

\noindent
\begin{minipage}[t]{0.48\linewidth}
\tzpsection{Wolfram Form}
\begingroup
\small
\[
\begin{aligned}
\mathfrak{W}_{\mathrm{TZP}}\!\left(\mathcal{K}_{1071}\right)
&\longmapsto
\left\langle \Omega^{\mathbb{T}},\; \Lambda,\; \Psi_{\mathrm{T}} \right\rangle
\end{aligned}
\]
\[
\mathcal{K}_{1071}
:=
\operatorname{HoldForm}\!\left(\mathcal{T}_{1071}\right)
\]
\[
\begin{aligned}
\operatorname{Admissible}_{\mathrm{TZP}}\!\left(\mathcal{K}_{1071}\right)
&\Longleftrightarrow
\mathfrak{W}_{\mathrm{TZP}}\!\left(\mathcal{K}_{1071}\right)\not\equiv\varnothing
\end{aligned}
\]
\textbf{Wolfram register:} canonical symbol lift\\
\textbf{Title lift:} MAGNETIC FIELD RAMP-UP\\
\textbf{Symbol key:} \(\Omega^{\mathbb{T}}\): Trawinistic boundary curvature\\ \(\Lambda\): Elsasser balance number\\ \(\Psi_{\mathrm{T}}\): Trawinistic wavefunction
\par
\endgroup
\end{minipage}\hfill
\begin{minipage}[t]{0.48\linewidth}
\tzpsection{TRAWIN Normalization}
\[
\tzpnormalizewrap{B(t) = \in t_0^t \mathrm{dot}{B} {ramp} dt \implies B {target} \approx 0.1 \text{ Tesla},\,- Magnetic field B(r,\theta,\phi) =  abla \times A,\,- Magnetic field: |B| > B_{\mathrm{max}} \to demagnetization risk}
\]

\small
\textbf{T}: Scaling Law is localized within the category theory and proof structure arc of the directory.\\
\textbf{R}: Geometric Constraint preserves the operative relation that magnetic field ramp-up requires.\\
\textbf{A}: Transfer Map fixes the admissible closure condition for the entry.\\
\textbf{W}: MAGNETIC FIELD RAMP-UP is rendered as a reusable operator-level statement.\\
\textbf{I}: The informational content of magnetic field ramp-up is retained rather than flattened.\\
\textbf{N}: MAGNETIC FIELD RAMP-UP closes as a distinct normalized TRAWIN object.
\normalsize
\end{minipage}

\tzpsection{Derivable Consequences}
The entry now reads through actual sourced equations, so its local arc is carried by explicit mathematical relations rather than a uniform quaternary certificate record shell.
\clearpage

\tzpentryheader{1071}{Oxford Dictionary of the Queens, NY English}

\begingroup\small
\tzpsection{Definition}
MAGNETIC FIELD RAMP-UP is a registry definition in Category Theory and Proof Structure with secondary pressure from Signal.

% BEGIN TZP CONTEXTUAL MEANING
\tzpsection{Contextual Meaning}
\begingroup\footnotesize\setlength{\emergencystretch}{3em}\sloppy
\textbf{Formal statement:} B(t) = int sub 0 to the power t dot B sub ramp dt implies B sub target approximately 0.1 text Tesla. \textbf{Contextual meaning:} The entry reads MAGNETIC FIELD RAMP-UP as a controlled technical headword whose equation layer, proof route, and registry role are interpreted together. \textbf{Derivable consequences:} The entry now reads through actual sourced equations, so its local arc is carried by explicit mathematical relations rather than a uniform quaternary certificate record shell. \textbf{Lemma support:} Scaling Law; Geometric Constraint; Transfer Map.\par
\endgroup
% END TZP CONTEXTUAL MEANING

\tzpsection{Technical Interpretation}
Technically, MAGNETIC FIELD RAMP-UP is read through the local equation chain and the TRAWIN normalization recorded on the entry page.

\tzpsection{Equation Grounding}
Equation anchors: B(t) = in t sub 0 to the power t dot B ramp dt implies B target approximately 0.1 Tesla; - Magnetic field B(r,theta,phi) = abla times A; and - Magnetic field: |B| > B sub max to demagnetization risk. Proof route: show that the stated equations form a coherent progression for the entry and remain compatible under the TRAWIN ordering. Formalization route: prove that the final closure relation follows from the preceding entry-specific equations.

\tzpsection{Interpretive Role}
The entry now reads through actual sourced equations, so its local arc is carried by explicit mathematical relations rather than a uniform quaternary certificate record shell.
\endgroup

\tzpsection{TZP Thesaurus}
\begin{enumerate}[leftmargin=1.6em,itemsep=6pt,topsep=4pt]
\item \textbf{Magnetized Carrier Domain Ramp Up}: entry-specific alternate wording for indexing, related literature, and local formal context.
\item \textbf{Magnetized Carrier Domain Rotating-Field Analogue}: rotating-field vocabulary for spin, handedness, inertial coupling, or magnetic transport.
\item \textbf{Carrier Domain Terminology}: entry-specific alternate wording for indexing, related literature, and local formal context.
\end{enumerate}

\tzpsection{Encyclopedia Mathematica}
\begin{samepage}
\noindent\begin{minipage}[t]{0.45\linewidth}
\vspace{0pt}
\raggedright
\scriptsize\textbf{Image reading.} The rendered manifold at right is the per-ID light plate for MAGNETIC FIELD RAMP-UP. It should be read as the visual companion to the entry's equation layer, not as a repeated template diagram.\par
\end{minipage}\hfill
\begin{minipage}[t]{0.53\linewidth}
\vspace*{-0.10in}
\raggedright
\tzpencyclopediaimage{1071}
\end{minipage}
\end{samepage}

\clearpage

\tzpentryheader{1071}{Direct Equation Progression and Proof Trajectory}

\tzpsection{Direct Equation Progression}
\begin{enumerate}[leftmargin=1.6em,itemsep=9pt,topsep=6pt]
\item \textbf{Scaling Law}
\[
B(t) = \in t_0^t \mathrm{dot}{B} {ramp} dt \implies B {target} \approx 0.1 \text{ Tesla}
\]
This fixes the first explicit relation carried by the entry rather than a generic certificate record normalization.

\item \textbf{Geometric Constraint}
\[
- Magnetic field B(r,\theta,\phi) =  abla \times A
\]
This adds the next operational equation that advances the entry beyond its opening definition.

\item \textbf{Transfer Map}
\[
- Magnetic field: |B| > B_{\mathrm{max}} \to demagnetization risk
\]
This closes the progression with an actual admissibility or closure relation instead of recycling the normalization shell.
\end{enumerate}

\tzpsection{Proof}
\textbf{Proof route.} show that the stated equations form a coherent progression for the entry and remain compatible under the TRAWIN ordering.

\textbf{Formal method.} prove that the final closure relation follows from the preceding entry-specific equations.

\medskip
\tzpproofartifacts{Isabelle audit: corpus classification recorded in appendix.}{Lean audit: compiled corpus certificate.}{Rocq audit: compiled corpus certificate.}

\tzpsection{Dependencies and Test Handle}
\begin{tzpsnapshot}
\tzpsnaprow{Dependencies}{\tzpidref{1070}, \tzpidref{1000}}
\tzpsnaprow{Node}{\tzpcode{registry_id_1071}}
\tzpsnaprow{Support Titles}{Scaling Law; Geometric Constraint; Transfer Map}
\tzpsnaprow{Test Handle}{If the cited entry equations cannot be made mutually admissible in one TRAWIN-normalized frame, the entry fails.}
\end{tzpsnapshot}

\tzpsection{Canonical Equation Links}
\tzpequationlinks
  {B(t) = \cdots}
  {- Magnetic field B(r,\theta,\phi) =  abla \times A}
  {|B| > B_{\mathrm{max}} \to demagnetization risk}
